\section*{Chapitre \ref{ch:m3:physical-commodity-markets} --- Marchés physiques de matières premières}

\begin{solution}{exo:m3:physical-commodity-markets:1}
Les cinq évaluations totalisent 403.25, soit une moyenne de 80.65~; en ajoutant 0.35, le prix est
\qty{81.00}{\$/\barrel}.
\end{solution}

\begin{solution}{exo:m3:physical-commodity-markets:2}
$78.40 + 1.35 + 0.04 = \qty{79.79}{\$/\barrel}$ CIF. Dans les deux cas, le risque
passe à l'acheteur lorsque le pétrole est à bord au port de chargement~: l'acheteur supporte
le risque de transport dans les deux cas~; avec le CIF, le vendeur paie le fret et
l'assurance qui couvre l'acheteur.
\end{solution}

\begin{solution}{exo:m3:physical-commodity-markets:3}
$84.10 \times 0.998 - 2.20 - 0.05 = \qty{81.68}{\$/\barrel}$.
\end{solution}

\begin{solution}{exo:m3:physical-commodity-markets:4}
Il achète 700 contrats à terme aujourd'hui et en vend 140 à chaque jour de fixation. Si un jour est écarté, le
prix est la moyenne de quatre évaluations~: il en vend 175 chacun des quatre jours.
\end{solution}

\begin{solution}{exo:m3:physical-commodity-markets:5}
\qty{1.08}{\$/\barrel}~: le coût couvert vaut $F_0$ plus la base moyenne plus le
différentiel (\cref{prop:m3:physical-commodity-markets:hedged}), dont l'écart-type
est proportionnelle à la volatilité de la base. Les opérations sur contrats à terme sont achetées et
revendues les mêmes jours que les prix physiques retenus, de sorte qu'elles se compensent exactement.
\end{solution}

\begin{solution}{exo:m3:physical-commodity-markets:6}
Les \omterm{def:m3:physical-commodity-markets:netback}{netbacks} valent 80.17 (Rotterdam) et 80.02 (Singapour)~: le prix livré à Singapour
doit augmenter de plus de \qty{0.15}{\$/\barrel}, pour dépasser
\qty{83.25}{\$/\barrel}.
\end{solution}

\begin{solution}{exo:m3:physical-commodity-markets:7}
Septembre 2011, $-\qty{27.31}{\$/\barrel}$~; le WTI a été en moyenne au-dessus du Brent dans 4 des
mois de janvier 2011 à août 2026 (115 sur les 320 mois depuis 2000, presque tous
avant 2011).
\end{solution}

\begin{solution}{exo:m3:physical-commodity-markets:8}
Le coût de la cargaison commence à flotter à la signature du contrat, et non à la date du \omterm{def:m3:physical-commodity-markets:bl}{connaissement} (B/L),
et il se fixe à un cinquième par jour sur les cinq jours de fixation, dont deux précèdent
la date du B/L. Acheter les 700 contrats à terme lorsque l'exposition naît, et en vendre 140 à chaque
jour de fixation au règlement-livraison du jour~; les vendre toutes «~une fois le prix connu~»
laisse les mouvements à terme des derniers jours sans aucun mouvement physique correspondant.
\end{solution}

\begin{solution}{pb:m3:physical-commodity-markets:1}
\textbf{1.} Zéro~: les deux prix flottent avec le Dated. \textbf{2.} L'exposition augmente de 140\,000
barils par jour à mesure que le prix d'achat se fixe, pour atteindre 700\,000 barils longs au jour 12. \textbf{3.}
700\,000 barils longs, du jour 12 au jour 27~: 16 jours. \textbf{4.} L'exposition baisse de
140\,000 par jour jusqu'à zéro au jour 32. \textbf{5.} Un trapèze
(la \cref{fig:m3:physical-commodity-markets:exposure}).
\textbf{6.} Vendre 140 contrats à terme par jour. \textbf{7.} Racheter 140 par jour.
\textbf{8.} Vente à découvert de 700. \textbf{9.} Un achat à prix flottant face à une vente à prix flottant ne
comporte aucune exposition au prix à plat. \textbf{10.} Un swap sur la moyenne Dated de chaque
\omterm{def:m3:physical-commodity-markets:formula}{période de prix} (un swap entre le Dated et les contrats à terme, du \cref{ch:m3:crude-oil}), ou les contrats à terme plus la
base résiduelle.
\textbf{11.} $1.10 - 0.20 - 0.65 - 0.03 = \qty{0.22}{\$/\barrel}$. \textbf{12.}
154\,000~\$. \textbf{13.} Vingt jours ouvrés entre les \omterm{def:m3:physical-commodity-markets:formula}{périodes de prix}~:
$1.80\sqrt{20} \times 700\,000 \approx$ 5.63~millions de dollars. \textbf{14.} La base
entre le Dated et les contrats à terme à chaque jour de fixation, la tolérance de volume, les surestaries et
les retards, le crédit de la contrepartie et les pertes opérationnelles. \textbf{15.} Les prix de vente
des jours 33--37~: l'exposition dure 21 jours à son maximum et les rachats interviennent cinq
jours plus tard~; la couverture doit rester ouverte et les contrats à terme pourraient devoir être reconduits.
\textbf{16.} Pour capter le fret et la prime de destination, et pour maîtriser le
transport. \textbf{17.} Le prix d'achat, payé environ trente jours après le chargement,
jusqu'à l'arrivée du produit de la vente, plus la marge sur les contrats à terme. \textbf{18.} Son P\&L
est le différentiel entre deux prix, et non le niveau des prix. \textbf{19.}
\emph{Résultat à retenir~:} 154\,000~\$ verrouillés par cargaison (\qty{0.22}{\$/\barrel}) face à une
écart-type non couvert d'environ 5.63~millions de dollars. \textbf{20.} Le spread
entre la valeur d'une matière première en deux lieux, deux dates ou deux formes, diminuée du coût de
son acheminement, et non le niveau des prix.
\end{solution}

\begin{solution}{iq:m3:physical-commodity-markets:1}
\omterm{def:m3:physical-commodity-markets:incoterms}{Franco à bord} (FOB)~: le vendeur charge sur le navire de l'acheteur et l'acheteur paie le fret et
l'assurance. \omterm{def:m3:physical-commodity-markets:incoterms}{Coût, assurance et fret} (CIF)~: le vendeur paie le fret et l'assurance jusqu'à destination. Dans les deux cas, le
risque passe à bord au port de chargement, de sorte que l'acheteur supporte le risque de transport.
\iqlookfor{la distinction entre qui paie et qui supporte le risque.}
\end{solution}

\begin{solution}{iq:m3:physical-commodity-markets:2}
Une moyenne se manipule plus difficilement par une seule transaction ou une seule mauvaise journée, étale la fixation des prix sur
la période pendant laquelle la cargaison est effectivement chargée ou livrée, et permet aux deux parties de se couvrir jour après
jour sur le \omterm{def:m3:physical-commodity-markets:layers}{marché papier}, liquide.
\iqlookfor{résistance à la manipulation, adéquation à la fenêtre opérationnelle, possibilité de couverture.}
\end{solution}

\begin{solution}{iq:m3:physical-commodity-markets:3}
À la date de la transaction, couvrir tout le volume, car l'exposition existe dès que
le prix flottera~; fermer ensuite $1/n$ de la couverture à chacun des $n$ jours de fixation
(les jours ouvrés du mois). Un swap sur moyenne mensuelle fait cela en un seul instrument.
\iqlookfor{exposition créée à la signature, dénouée aux jours de fixation, le swap comme couverture naturelle.}
\end{solution}

\begin{solution}{iq:m3:physical-commodity-markets:4}
Le règlement-livraison est calculé par une règle à partir des transactions en bourse et il est reproductible~;
l'évaluation relève d'un jugement d'une équipe éditoriale portant sur les cours acheteurs, les cours vendeurs et les transactions sur une
fenêtre, parfois avec peu de transactions, et peut être révisée. Sa méthodologie, sa fenêtre et
son horodatage comptent, et son historique est une donnée sous licence.
\iqlookfor{reproductibilité, risque méthodologique, licences.}
\end{solution}

\begin{solution}{iq:m3:physical-commodity-markets:5}
Sur la base entre l'indice de référence couvert et la formule de la cargaison, sur les décalages de calendrier
(retards, tolérances), sur le défaut de la contrepartie, sur les pertes opérationnelles et les surestaries,
et sur le financement des appels de marge sur ses couvertures tant que le gain physique n'est
pas encore réalisé.
\iqlookfor{base, opérations, crédit et liquidité plutôt que prix à plat.}
\end{solution}

\begin{solution}{iq:m3:physical-commodity-markets:6}
Saisie des transactions avec la formule, l'indice de référence et les jours de fixation de chaque jambe~; un calendrier
des jours ouvrés par indice de référence~; un traitement nocturne qui calcule la fraction fixée de chaque
jambe, agrège l'exposition par indice de référence et par jour, la compare aux contrats à terme et swaps
détenus, et produit les transactions à solder~; alertes sur les évaluations manquantes et sur les jambes
dont les dates se déplacent.
\iqlookfor{le modèle de données jambe et jour de fixation, l'agrégation par indice de référence, le rapprochement avec les couvertures.}
\end{solution}
